\section{Introduction}

New-onset postoperative atrial fibrillation (POAF) complicates roughly one in four isolated coronary artery bypass grafting (CABG) procedures, driving stroke, heart failure, prolonged intensive care stay, readmission, and excess long-term mortality.\cite{aranki1996predictors,villareal2004,elchami2010poaf,lapar2014poaf,dobrev2019postoperative,jawitz2020poaf,caldonazo2023meta} Because the most effective preventive strategies---including intraoperative posterior left pericardiotomy and targeted perioperative prophylaxis---must be selected \emph{before skin incision},\cite{gaudino2021palacs,chatterjee2026sts} clinicians require pre-incision risk models that are accurate, well calibrated, and fast enough for real-time hospital deployment.

Despite dozens of published POAF prediction models, critical appraisals under PROBAST standards show that apparent derivation performance (median ROC-AUC 0.71) routinely degrades on external validation (median 0.61), with more than 80\% rated at high risk of bias.\cite{fields2025poaf,lopeztuesta2026poaf,pandey2023poaf,wolff2019probast,moons2025probastai} Two structural evaluation flaws obscure real-world clinical performance in national surgical registries. First, institutional POAF incidence varies threefold across U.S.\ cardiac centers due to differences in case mix, perioperative protocols, and rhythm ascertainment.\cite{filardo2017telemetry,perezgrovas2023jaha} Consequently, naive calendar holdouts allow nearly all test hospitals to appear in training, conflating genuine patient-level generalization with hospital familiarity.\cite{steyerberg2016prediction,riley2019external} Second, modern tabular foundation models (TFMs) and multi-layer automated machine learning (AutoML) pipelines have transformed structured-data benchmarks,\cite{hollmann2025tabpfn,erickson2025tabarena,erickson2020autogluon,jaeger2026tabpfn35,qu2026kumotabular,kong2026tabfm} yet clinical evaluations rarely test them side-by-side across both \emph{strictly unseen hospitals} (cold-start deployment) and \emph{seen hospitals} (in-network deployment), nor do they audit the severe inference latency and GPU memory costs imposed by in-context foundation models at hospital scale.

Using the full national Society of Thoracic Surgeons Adult Cardiac Surgery Database (STS ACSD) v2.81 isolated-CABG registry (368,132 procedures across 1,091 U.S.\ centers),\cite{sts2014v281,shahian2018sts} we establish a leakage-audited pre-incision benchmark across both validation regimes. On identical 50,000-patient training cohorts and strictly 207 pre-incision variables, we compare five leading binary tabular AI alternatives---Google TabFM v1.0.0, NVIDIA KumoTabular-Large v2, PriorLabs TabPFN-3.5, tuned CatBoost, and AutoGluon-Tabular 1.5 (\cli{best\_quality})---directly against NEVERHUMAN, our proprietary pre-incision risk model, evaluating discrimination, calibration, clinical net benefit, subgroup equity, and end-to-end training and inference speed.
